\documentclass[11pt]{article}
\usepackage[T1]{fontenc}
\usepackage[utf8]{inputenc}
\usepackage{lmodern}
\usepackage[margin=1in]{geometry}
\usepackage{amsmath,amssymb,amsthm,mathtools}
\usepackage{microtype}
\usepackage{booktabs}
\usepackage[hidelinks]{hyperref}
\newtheorem{theorem}{Theorem}[section]
\newtheorem{proposition}[theorem]{Proposition}
\newtheorem{corollary}[theorem]{Corollary}
\newcommand{\at}[1]{\langle #1\rangle}
\newcommand{\ch}[1]{\mathbin{\oplus_{#1}}}
\newcommand{\ow}[1]{\mathbin{\star_{#1}}}
\newcommand{\Pow}{\mathcal P_{+}}
\newcommand{\B}{\mathsf B}
\title{Small Owned Choice:\\Neutral Union, Occurrence Bounds, and Unary Contexts}
\author{KFOL research working note}
\date{28 October 2026}
\begin{document}
\maketitle
\begin{abstract}
Fix a finite set of roles. Consider finite public games with only
role-owned binary choice and a unique winner at each leaf. Observe the
possible winners of pure subgame-perfect equilibria. This observation has
a small compositional algebra. Three distinct owners define ordinary set
union, using two occurrences of each argument. We prove that every union
term needs at least two occurrences of each argument, even with arbitrary
ground constants. We classify all unary term functions and derive their
exact number. The union construction also works with identical component
strategies in its repeated copies, under ordinary unrestricted deviations.
The game interpretation, union identity, occurrence obstruction, uniform-copy
result, and upper half of the classification are kernel-checked in Lean.
The remaining arguments are given explicitly below. This note records a
candidate algebraic contribution; priority has not been established.
\end{abstract}

\section{Why start here?}

The motivating question is whether equally primitive player-owned choices
can support a useful calculus. A small answer should first explain what
composition preserves, which equations it validates, and what an observer
can distinguish. Dependent types, control, hidden information, and beliefs
should enter when a specific theorem requires them.

This note fixes one particularly small experiment:
\[
  G ::= \at i \mid x \mid G\ch r G, \qquad i,r\in I.
\]
The set $I$ has $K$ roles. A closed term is a finite, perfect-information
tree: at a node owned by $r$, role $r$ chooses one child. Leaf $\at i$
gives utility $1$ to $i$ and $0$ to every other role. Histories are public;
there is no chance, recursion, hidden state, or sharing of decision nodes.
Substituting a term at two occurrences produces two subtrees.

The observation $\B(G)$ is the set of winners attained by full pure
subgame-perfect equilibrium (SPE) profiles. It is a set of possible
equilibrium outcomes, not a mixed strategy or an individual guarantee.
This is a semantic choice to investigate, rather than a claim that rational
equilibrium is the only reading of an unbiased choice quantifier.

Owner-indexed game syntax is established background
\cite{Abramsky2007,OldeLoohuisVenema2010}. Full contingent strategies and
dependent game trees also have established type-theoretic treatments
\cite{EscardoOliva2023}. We use these foundations directly. The proposed
contribution lies in the algebra and context classification under the
specified observation.

\section{From full strategies to a finite algebra}

A full profile gives an action at \emph{every} internal node, including
nodes avoided by that profile. At a node it therefore consists of a root
action and a full profile for each child. A unilateral deviation by $i$
may replace all actions owned by $i$, in both child trees. It leaves every
other player's actions fixed. A profile is Nash if no such deviation
improves the deviator's utility. It is SPE if its restriction to every
subtree is Nash.

For a fixed profile, let $a,b$ be the winners prescribed in its two child
subtrees. Assuming both child profiles are SPE, the parent profile is SPE
exactly when it chooses $r$ whenever $a=r$ or $b=r$. Otherwise either
child is optimal for $r$. This characterization follows from whole-player
deviations: a deviation that wins in a child would already improve that
child's profile. It is not an assumption that checking only a current
action is generally sufficient in games of partial information.

For nonempty subsets $S,T\subseteq I$, put
\begin{equation}\label{eq:own}
 S\ow r T =
 \begin{cases}
 \{r\},&S=\{r\}\text{ or }T=\{r\},\\
 S\cup T,&\text{otherwise}.
 \end{cases}
\end{equation}

\begin{theorem}[Exact public-tree interpretation]\label{thm:adequacy}
Every closed tree has a full pure SPE profile, and
\[
 \B(\at i)=\{i\},\qquad
 \B(G\ch r H)=\B(G)\ow r\B(H).
\]
Every element of the right-hand side has a full SPE witness on the left.
\end{theorem}
\begin{proof}
Existence follows inductively by choosing an owner-winning child whenever
one is available in the chosen child profiles. If one child has only
$r$-winning SPE profiles, every parent SPE must win for $r$, and one
exists. Otherwise each child has an SPE profile losing for $r$. For any
SPE winner $i$ in the left child, combine its witness with a losing-for-$r$
SPE in the right child and choose left. This is locally optimal, whether
or not $i=r$. The symmetric construction covers the right child.
Conversely, a parent SPE selects the winner of a child SPE. These facts
give both inclusions in~\eqref{eq:own}.
\end{proof}

This is a backward-induction result for a restricted utility model. Its
role here is to justify the finite algebra $\Pow(I)$, not to claim new
game theory. \texttt{Basic.lean} defines full profiles, whole-player
deviations, Nash equilibrium, and SPE separately, then proves the theorem
as \texttt{winners\_node}.

Each $\star_r$ is associative, commutative, and idempotent. For
associativity, a three-input product is $\{r\}$ if any input is $\{r\}$,
and is their union otherwise. Nonemptiness ensures that a union of inputs
different from $\{r\}$ cannot become $\{r\}$. Role permutations transport
these operations and the leaf constants.

For two roles, the ground-generated carrier consists only of the two
singletons and recovers Boolean conjunction/disjunction after choosing a
role as truth. With three roles there is genuine ambiguity, such as
$\{0,1\}$, and the next construction generates every nonempty subset.

\section{Neutral union and its minimum occurrence cost}

\begin{theorem}[Neutral union]\label{thm:union}
For pairwise distinct $r,s,t$,
\begin{equation}\label{eq:union}
 J(x,y)=(x\ow r y)\ow t(x\ow s y)=x\cup y.
\end{equation}
Consequently, if $K\geq3$, every nonempty subset of $I$ is the
denotation of a closed owned-choice tree.
\end{theorem}
\begin{proof}
If neither inner operation collapses, both return $x\cup y$. If just one
collapses, its owner singleton belongs to the other result, and the outer
owner is different. If both collapse, the inputs must be the two distinct
owner singletons; the outer operation returns their union. Folding $J$
over the leaf singletons realizes every nonempty subset.
\end{proof}

The definition uses three owned-choice nodes and two occurrences of each
input. Although it names three owners, its denotation is independent of
the distinct triple chosen. It is an equality of equilibrium-winner
observations, not an isomorphism of trees or strategy spaces.

\begin{theorem}[Occurrence lower bound]\label{thm:lower}
Suppose $K\geq3$. In the syntax of owned binary choices and arbitrary
nonempty ground constants, any term that denotes $x\cup y$ for every
$x,y\in\Pow(I)$ contains at least two occurrences of $x$ and at least two
occurrences of $y$. Hence three binary nodes are optimal.
\end{theorem}
\begin{proof}
Zero occurrences cannot express dependence on that input. Suppose $x$
occurs once. Read the unique path from its occurrence to the root, using
commutativity to put the current value first. It has the form
\[
 X\longmapsto X\ow{r_1}g_1(Y)
   \longmapsto\cdots\longmapsto (\cdots)\ow{r_n}g_n(Y),
\]
where each $g_j(Y)$ is nonempty and independent of $X$. We allow arbitrary
such functions, which only strengthens the obstruction.

For fixed $Y$, if some $g_j(Y)=\{r_j\}$, that step erases all dependence
on $X$, and later steps cannot recover it. For $Y=\{r\}$ this is
incompatible with $X\cup\{r\}$, which varies with $X$. Therefore no such
side value equals its owner singleton. Every step then preserves all
members of its current input: if the input is the owner singleton it
remains so, and otherwise the operation takes a union.

If the path is empty, its value is $X$, not union. Otherwise let $r$ be
its innermost owner and set $Y=\{r\}$. Its first side value differs from
$\{r\}$ and is nonempty, so contains some $k\ne r$. Choose
$j\notin\{r,k\}$. Starting with $X=\{j\}$, the first step introduces
$k$. No later step can remove it. The required result $\{j,r\}$ excludes
$k$, a contradiction. Interchange the variables for the other bound.
\end{proof}

The path obstruction is \texttt{no\_single\_use\_union} in Lean. The
syntax-to-path correspondence is the direct decomposition of a term at
its unique input occurrence, as used in the proof.

The cost measured here is \emph{syntactic occurrence}. A play enters just
one of the four component occurrences. The theorem establishes neither
runtime duplication nor a need for a contraction rule in a particular
linear type system: additive branching may use a context in both
exclusive branches.

\section{Repeated descriptions with uniform component strategies}

Repeated occurrences normally have separately chosen full strategy
tables. Fixing one component SPE profile in every copy gives a different
observation in general.

\begin{proposition}[Uniform-profile observation]
Let $f[x]$ be any unary owned-choice context and $G$ any closed component.
The winners of ordinary SPE profiles that use one and the same full SPE
profile of $G$ in every copy are exactly
\[
 U(f)(\B(G))=\bigcup_{g\in\B(G)}f(\{g\}).
\]
Thus uniform copies preserve the ordinary observation at $S$ precisely
when $f(S)=\bigcup_{g\in S}f(\{g\})$.
\end{proposition}
\begin{proof}
Fix the component profile. Replace every copy with the leaf bearing its
winner. Local optimality inside the component is already supplied by
that profile; local optimality at each outer node depends only on child
winners. Induction over the context therefore preserves SPE and the
selected winner. Take the union over component SPE profiles.
\end{proof}

This generalization was identified in the external review
\path{notes/review-results-2026-10-28.md}. The next theorem is its
specific positive case for neutral union, with a generic Lean proof.

\begin{theorem}[Uniform copies suffice]\label{thm:uniform}
Fix a full SPE profile $p$ of $G$ and a full SPE profile $q$ of $H$, with
winners $g,h$. In
\[
  (G\ch r H)\ch t(G\ch s H),
\]
use exactly $p$ in both copies of $G$ and exactly $q$ in both copies of
$H$. As the three outer actions vary, the winners of ordinary SPE
profiles are exactly $\{g,h\}$. Therefore the winners of ordinary SPE
profiles satisfying this copy-equality condition are
$\B(G)\cup\B(H)$.
\end{theorem}
\begin{proof}
For any fixed three outer actions, local optimality at those nodes depends
only on the child winners $g,h$. Local optimality inside a component is
already supplied by $p,q$. Thus the expanded full profile is SPE exactly
when the corresponding profile of
$(\at g\ch r\at h)\ch t(\at g\ch s\at h)$ is SPE. By
Theorems~\ref{thm:adequacy} and~\ref{thm:union}, its possible winners are
$\{g,h\}$. Conversely, a uniform ordinary SPE restricts to SPE profiles
in its component copies. Taking the union over $p,q$ proves the claim.
\end{proof}

This is \texttt{uniform\_copies\_winners} in Lean. Importantly, we
\emph{select} profiles with equal component copies from the ordinary SPE
profiles. We do not restrict deviations to preserve that equality. The
result therefore survives this test of reuse without changing the
equilibrium concept. It does not identify outer decisions with component
decisions, bind a previously realized value, or introduce an information
set spanning two histories. Those are separate interface experiments.

The term $D$ below supplies a separating case: it fixes every singleton
but adds winners on some ambiguous inputs, so $U(D)$ is the identity
while $D$ is not. Uniform observation is also not preserved by arbitrary
substitution of an expression for the input variable. The follow-up note
\texttt{structure.tex} gives an explicit substitution counterexample and
develops the multivariable algebra. The theorem above makes no general
claim that repetition can be treated as shared selection.

\section{All unary contexts}

Assume $K\geq3$. A unary term function uses one variable, any number of
its occurrences, owned choices, and leaf constants. By
Theorem~\ref{thm:union}, arbitrary nonempty ground constants are available.
Write $\mathsf{Reg}(C,E)$ for the following function, with parameters
\[
 C\subseteq I,\qquad E_i\subseteq C\setminus\{i\}\quad(i\in I):
\]
\begin{equation}\label{eq:normal}
 \mathsf{Reg}(C,E)(S)=
 \begin{cases}
 \{i\}\cup E_i,&S=\{i\},\\
 S\cup C,&|S|\geq2.
 \end{cases}
\end{equation}

\begin{theorem}[Unary normal forms]\label{thm:unary}
The unary term functions are exactly the nonempty constant functions and
the functions~\eqref{eq:normal}. The regular parameters $C,(E_i)$ are
unique. The only overlap with the constant functions is the constant
function $I$.
\end{theorem}
\begin{proof}
\emph{Upper bound.} The variable is regular with $C=\varnothing$ and
empty rows, and each constant is constant. Suppose two child functions
are regular with extra sets $C,D$. On a multiwinner input $S$, each child
contains $S$ and cannot equal an owner singleton. Their owned combination
is their union, $S\cup C\cup D$. On singleton $\{i\}$, both contain $i$;
their owned combination still contains $i$ and is contained in
$\{i\}\cup C\cup D$. This is regular.

If one child is constant $D$ and the other regular with extra set $C$,
then $D=\{r\}$ makes the result constant. Otherwise the regular child's
multiwinner values cannot collapse and the result is $S\cup C\cup D$.
On an atomic input, the operation preserves $i$ and stays inside
$\{i\}\cup C\cup D$. Two constants give a constant. This proves closure
and the upper classification by induction. The Lean theorem
\texttt{unary\_normal\_form\_upper} formalizes this argument, including
extensivity of the regular functions.

\emph{Synthesis.} For every $k\in C$, prescribe the allowed singleton
input roles
\[
 A_k=\{i\ne k:k\in E_i\}.
\]
Start with $g_k(x)=J(x,\at k)$. For each $i\ne k$ outside $A_k$, replace
$g_k(x)$ by $g_k(x)\ow i x$. At singleton input $i$, this suppresses the
extra winner $k$ by returning $\{i\}$. At another singleton, it preserves
the current value, which contains that other role. On a multiwinner
input $S$, both operands contain $S$, so the result remains $S\cup\{k\}$.
The order of these suppressions is irrelevant. Join $x$ with all the
$g_k(x)$ using $J$. This realizes precisely~\eqref{eq:normal}; for empty
$C$, use $x$.

\emph{Uniqueness.} For any $k$, choose a two-element $S$ omitting $k$.
Then $k\in C$ exactly when $k\in\mathsf{Reg}(C,E)(S)$, so $C$ is uniquely
determined. Each $E_i$ is the singleton-input result minus $i$. A regular
function is extensive. If it is constant, its value must contain every
singleton and hence equal $I$. That constant is regular with $C=I$ and
$E_i=I\setminus\{i\}$.
\end{proof}

\begin{corollary}[Exact count and a small test set]
There are exactly
\begin{equation}\label{eq:count}
  (1+2^{K-1})^K+2^K-2
\end{equation}
unary term functions. Equality of two unary functions is determined by
their values on singleton and two-element input sets.
\end{corollary}
\begin{proof}
For fixed $C$ there are $2^{|C|(K-1)}$ row families. Summing over $C$
gives $(1+2^{K-1})^K$. Add the $2^K-1$ constants and subtract the single
overlap. Singleton/pair tests determine all regular parameters. They
also distinguish constants: a constant agreeing with a regular function
on every singleton must contain every role, hence be $I$; pair tests
then force $C=I$ and all rows to be full.
\end{proof}

Thus there are $131$ unary functions for $K=3$, $6575$ for $K=4$, and
$1419887$ for $K=5$. The first number is independently recovered by
exhaustive algebraic closure in \texttt{check.py}. The script constructs
every parameterized function by actual owned-choice operations for
$K=3,4$ and checks separation by singleton/pair tests. It does not
enumerate all terms for $K=4,5$; the general classification supplies the
upper bound.

The parameters are also closed under context composition. If the outer
function has parameters $(C,E)$ and the inner $(D,F)$, their composite is
regular with extra set $C\cup D$ and rows
\[
 H_i=\begin{cases}
 E_i,&F_i=\varnothing,\\
 (F_i\cup C)\setminus\{i\},&F_i\ne\varnothing.
 \end{cases}
\]
If either function is constant, evaluate the appropriate constant case.
Indeed, a multiwinner input stays multiwinner through both functions. At
input $\{i\}$, the inner result is still a singleton exactly when
$F_i=\varnothing$; otherwise the outer function adds $C$.
This supplies an explicit finite monoid of unary contexts rather than a
list of isolated identities.

Atomic inputs alone do not determine a context. For $I=\{a,b,c\}$, the term
\[
 D(x)=x\ow c\bigl((x\ow a\{b\})\ow b\{a\}\bigr)
\]
fixes every singleton but sends each multiwinner set $S$ to
$S\cup\{a,b\}$. It has only two occurrences of $x$. Thus a calculus that
checks equations only on the winner constants misses real distinctions.

\section{Which observations retain the information?}

\begin{proposition}[Unique-winner tests collapse ambiguity]\label{prop:obs}
Every unary context function is either constant or extensive. If its
observation is only whether its result equals a designated singleton,
all sets with at least two members are contextually indistinguishable.
If membership of a designated winner is observable, distinct sets are
already distinguished by the identity context.
\end{proposition}
\begin{proof}
The first statement follows either directly by induction (a constant
owner-singleton absorbs; otherwise a nonconstant child is preserved), or
from Theorem~\ref{thm:unary}. A constant context has the same result on all
inputs. An extensive context cannot turn a multiwinner input into a
singleton. Membership tests distinguish sets at an element of their
symmetric difference.
\end{proof}

The first two statements are proved generically in Lean as
\texttt{unary\_constant\_or\_extensive} and
\texttt{ambiguity\_unobservable\_by\_unique\_winner}; they do not depend
on enumerating the finite carrier. This settles the observation issue in
Target C of the realization-and-sharing proposal for the public unary
fragment. It also identifies a specific way an apparently richer logic
could accidentally lose its intended content.

For comparison, let $q(S)=\{i:S=\{i\}\}$. For nonempty operands,
\[
 q(S\ow r T)_r=q(S)_r\lor q(T)_r,\qquad
 q(S\ow r T)_i=q(S)_i\land q(T)_i\quad(i\ne r).
\]
These are precisely the coordinate operations in the guarantee algebra
of Olde Loohuis and Venema, Definition 3.20
\cite{OldeLoohuisVenema2010}, restricted to its empty and singleton
values. Induction also shows $\B(G)=\{i\}$ exactly when $i$ can force a
win in this finite unique-winner fragment. This comparison explains the
collapse; it is not the chosen meaning of $\B$.

The full algebra validates different equations. For example, with
$r=0,s=1$, $A=\{0,1\}$, $B=\{1\}$, and $C=\{2\}$,
\[
 A\ow0(B\ow1 C)=\{0,1\},\qquad
 (A\ow0 B)\ow1(A\ow0 C)=\{0,1,2\}.
\]
Thus the mutual distributivity axiom of their multiplayer algebras
fails here, even though every $\star_r$ distributes over the neutral
union operation. Reusing their complete equational calculus would
therefore change the semantics.

\section{The partial-information boundary and the next calculus}

Full sequential equilibrium (SE) is not a small primitive for this
experiment. It requires assessments containing strategies and beliefs,
sequential rationality at information sets, and a consistency condition.
The current results give no compression theorem for those data.

There is also a precise reason not to extend the public recurrence by
fiat. Suppose Alice chooses a bit and Bob guesses it, winning exactly
when he guesses correctly. With public observation, every pure SPE
winner is Bob. Hide the bit from Bob while keeping the same tree and
terminal winners. At a matching pair Alice can switch her bit; at a
mismatching pair Bob can switch his guess. Hence there is no pure Nash
equilibrium. Bob's information set crosses both branches, so neither
Bob node roots a proper continuation subgame. SPE supplies no additional
restriction beyond Nash here. It is the refinement that becomes vacuous,
not the definition of Nash equilibrium.

The public winner-set observation therefore cannot support arbitrary
hidden-information composition. Switching to SE changes the semantic
problem; it does not repair the existing recurrence for free. This is a
boundary example, not a general impossibility theorem for small calculi
with hidden information.

The next step should keep the public core and specify one shared binary
decision as an interface. Ask which equations survive and which pair of
components equal under $\B$ becomes distinguishable. The uniform-copy
theorem supplies a positive case: reusing the same full component profile
does not itself destroy neutral union. The hidden-bit example supplies a
negative case. A useful extension must characterize the gap, with
whole-player deviations preserved throughout.

For a first typed calculus, the immediate requirement is a substitution
theorem for a chosen interface class. Distinguishing a reusable
description, a previously chosen value, and a shared policy entry should
follow from that theorem. The occurrence bound alone does not determine
structural typing rules. A dependent calculus and a treatment of SE remain
possible later targets, rather than premises of the current algebra.

\section{Evidence, attribution, and remaining work}

\begin{center}
\small
\begin{tabular}{@{}p{0.46\linewidth}p{0.46\linewidth}@{}}
\toprule
Result & Present evidence\\
\midrule
Full-strategy SPE adequacy & Generic Lean proof, including whole-player deviations\\
Neutral union & Generic Lean proof\\
Single-use obstruction & Generic Lean path theorem; syntax decomposition above\\
Uniform copies & Generic Lean proof for each fixed pair of component SPE profiles\\
Unary classification upper bound & Generic Lean proof\\
Unary synthesis, uniqueness, count, composition & Mathematical arguments above; finite checks for $K=3,4$\\
Unique-winner observation collapse & Generic Lean proof\\
Originality of the algebraic package & Open: targeted literature comparison, not a priority proof\\
\bottomrule
\end{tabular}
\end{center}

The Lean development imports only \texttt{Std} and uses installed Lean
4.34.1. It has no admitted lemmas or added axioms. Printed dependencies
are Lean's standard classical foundations: propositional extensionality,
classical choice, and quotient soundness. The Python script provides
finite cross-checks; it is not a replacement for the generic proofs.

The closest comparisons separate syntax, guarantees, and rational
strategies. Abramsky supplies owner-indexed game constructions;
Olde Loohuis--Venema supplies a different, guarantee-based algebra;
Escard\'o--Oliva supplies full optimal-strategy machinery.
Bolt--Hedges--Zahn's obstruction concerns a specified product of
nondeterministic selection functions\cite{BoltHedgesZahn2019}. This note
uses neither that product nor a claim that arbitrary component winner
sets compose under partial information.

The union term was already present in
\texttt{research/proposal-realization-and-sharing.md}; it is not a new
discovery of this note. The results worth a priority investigation are
its optimal syntactic occurrence bound, the uniform-copy preservation
theorem, and the unary normal-form classification. Searches did not find
these exact statements, but that does not establish novelty. A complete
equational axiomatization and the first nontrivial shared-decision
interface remain open. This is a concrete algebraic research core, not
yet a complete new logic or a CIC-style type theory.

\bibliographystyle{plain}
\bibliography{references}
\end{document}
